\documentclass[a4paper]{article}

\newcommand{\notetitle}{Lecture 02: Reddit at Scale --- gobernanza de comunidades, API de la plataforma e interfaces de IA}
\newcommand{\noteauthors}{Elaborado a partir de la entrevista con Steve Huffman, los materiales del curso Stanford Winter 2025 y los subtítulos con marcas de tiempo}
\newcommand{\notedate}{Curso de invierno de 2025 · Versión reescrita en 2026}
\newcommand{\videochannel}{Publicación histórica del curso Stanford CS153 (actualmente en private)}
\newcommand{\videopublishdate}{2025-05-06}
\newcommand{\videoduration}{48:22}
\newcommand{\videourl}{https://www.youtube.com/watch?v=yeA-opPcYxk}
\newcommand{\videocoverpath}{cover.jpg}

% Espanol (Mexico) con XeLaTeX: fontspec para fuentes Unicode, polyglossia
% para la division silabica y los nombres del documento en la variante mexicana.
\usepackage{fontspec}
\usepackage{polyglossia}
\setdefaultlanguage[variant=mexican]{spanish}
% polyglossia no traduce los nombres de listings.
\AtBeginDocument{\providecommand{\lstlistingname}{}\renewcommand{\lstlistingname}{Listado}%
  \providecommand{\lstlistlistingname}{}\renewcommand{\lstlistlistingname}{Índice de listados}}
\usepackage{amsmath,amssymb}
\usepackage{graphicx}
\usepackage[margin=2.35cm]{geometry}
\usepackage[most]{tcolorbox}
\usepackage{etoolbox}
\usepackage{listings}
\usepackage{xcolor}
\usepackage{hyperref}
\usepackage{booktabs,longtable,tabularx,array}
\usepackage{subcaption}
\usepackage{float}
\usepackage{enumitem}
\usepackage{microtype}
\usepackage{fancyhdr}
\usepackage{needspace}

\definecolor{kbBack}{HTML}{F2F6FA}
\definecolor{kbFrame}{HTML}{9DB4CC}
\definecolor{kbTitle}{HTML}{52708F}
\definecolor{ibBack}{HTML}{FBF5E7}
\definecolor{ibFrame}{HTML}{C9A961}
\definecolor{ibTitle}{HTML}{7D5F1E}
\definecolor{wbBack}{HTML}{FCF2F0}
\definecolor{wbFrame}{HTML}{D6A096}
\definecolor{wbTitle}{HTML}{9E4F43}
\definecolor{dbBack}{HTML}{F4F7F3}
\definecolor{dbFrame}{HTML}{AEC2AC}
\definecolor{dbTitle}{HTML}{5F7F64}

\tcbset{softbox/.style={
  enhanced,breakable,boxrule=0.8pt,arc=2pt,left=3mm,right=3mm,
  top=2mm,bottom=2mm,coltitle=white,fonttitle=\bfseries,
  toptitle=0.6mm,bottomtitle=0.6mm
}}
\newtcolorbox{knowledgebox}[1]{softbox,colback=kbBack,colframe=kbFrame,colbacktitle=kbTitle,title={#1}}
\newtcolorbox{importantbox}[1]{softbox,colback=ibBack,colframe=ibFrame,colbacktitle=ibTitle,title={#1}}
\newtcolorbox{warningbox}[1]{softbox,colback=wbBack,colframe=wbFrame,colbacktitle=wbTitle,title={#1}}
\newtcolorbox{dialoguebox}[1]{softbox,colback=dbBack,colframe=dbFrame,colbacktitle=dbTitle,title={#1}}

\lstset{
  language=Python,basicstyle=\ttfamily\small,
  keywordstyle=\color{kbTitle}\bfseries,stringstyle=\color{wbTitle},
  commentstyle=\color{dbTitle},breaklines=true,frame=single,numbers=left,
  numberstyle=\tiny\color{gray},captionpos=b,columns=fullflexible,
  keepspaces=true,showstringspaces=false,extendedchars=true
}
\BeforeBeginEnvironment{lstlisting}{\Needspace{12\baselineskip}}

\hypersetup{colorlinks=true,linkcolor=kbTitle,urlcolor=kbTitle,citecolor=wbTitle}
\setlist{itemsep=0.25em,topsep=0.4em}
\setlength{\parindent}{2em}
\setlength{\parskip}{0.25em}
\newcolumntype{L}[1]{>{\raggedright\arraybackslash}p{#1}}

\providecommand{\notetitle}{Notas del curso: [tema del curso]}
\providecommand{\noteauthors}{Elaboradas a partir de los materiales públicos de Stanford CS153}
\providecommand{\notedate}{Primavera de 2026}
\providecommand{\videochannel}{Stanford Online}
\providecommand{\videopublishdate}{2026}
\providecommand{\videoduration}{[duración del video]}
\providecommand{\videourl}{}
\providecommand{\videocoverpath}{cover.jpg}
\providecommand{\coursesite}{https://cs153.stanford.edu/}
\providecommand{\repourl}{https://github.com/hqhq1025/ai-course-notes}
\providecommand{\skillversion}{wdkns/wdkns-skills@39f1a04}

\newcommand{\figsource}[1]{\par\vspace{-0.3em}{\footnotesize\color{gray}\emph{Fuente: #1}}\par}
\newcommand{\teachervoice}[1]{\begin{knowledgebox}{Nota de clase}#1\end{knowledgebox}}
\newcommand{\term}[1]{\textbf{#1}}

\pagestyle{fancy}
\fancyhf{}
\fancyfoot[C]{\thepage}
\fancyfoot[R]{\footnotesize\color{gray}\href{\repourl}{AI Course Notes}}
\renewcommand{\headrulewidth}{0pt}
\renewcommand{\footrulewidth}{0pt}

\newcommand{\makecscover}{
\begin{titlepage}
\centering
\vspace*{0.7cm}
{\Large Stanford CS153 · Frontier Systems\par}
\vspace{0.8cm}
{\huge\bfseries \notetitle\par}
\vspace{0.6cm}
{\large \noteauthors\par}
\vspace{0.25cm}
{\large \notedate\par}
\vspace{0.9cm}
\ifdefempty{\videocoverpath}{}{
  \includegraphics[width=0.86\textwidth,height=0.43\textheight,keepaspectratio]{\videocoverpath}\par
}
\vfill
\begin{tcolorbox}[width=0.92\textwidth,colback=black!2!white,colframe=black!45,boxrule=0.7pt,arc=2pt]
\textbf{Autor o canal del video}: \videochannel\par
\textbf{Fecha de publicación}: \videopublishdate\par
\textbf{Duración del video}: \videoduration\par
\textbf{Sitio del curso}: \href{\coursesite}{\nolinkurl{\coursesite}}\par
\ifdefempty{\videourl}{}{\textbf{Enlace del video}: \href{\videourl}{\nolinkurl{\videourl}}\par}
\textbf{Normas de elaboración}: \skillversion\ + las normas source-first / teacher-voice / visual-QA de este repositorio
\end{tcolorbox}
\end{titlepage}
\tableofcontents
\newpage
}

\graphicspath{{./}{images/}}
\setcounter{tocdepth}{1}

\begin{document}
\makecscover

\begin{warningbox}{Nota sobre la disponibilidad del material fuente}
La grabación de esta sesión que se había subido a YouTube cambió su visibilidad a private el 2026-08-11. Estas notas usan los subtítulos con marcas de tiempo y la miniatura oficial que se conservan en el repositorio, y redibujan con `lecture02-diagrams.py` 12 diagramas conceptuales didácticos. Cada mecanismo que muestran las figuras puede verificarse en el tramo correspondiente de los subtítulos. No son capturas de las slides del video original y no deben confundirse con páginas que haya mostrado el ponente.
\end{warningbox}

\section{De un ciclo simple a una plataforma de comunidades}

La historia de Reddit sirve para observar cómo un sistema de internet pasa de ser un prototipo de producto a convertirse en infraestructura social. El objetivo inicial no era «construir decenas de miles de comunidades», sino permitir que los usuarios descubrieran contenido nuevo en la web: enviar enlaces, votar para ordenarlos y discutir en torno a ellos. Este ciclo era lo bastante simple para que los fundadores lo implementaran en pocas semanas, y lo bastante abierto para que los usuarios transformaran poco a poco un agregador de noticias en una red de comunidades con identidades, normas y conocimientos distintos. Entender esta trayectoria de evolución explica mejor que memorizar el año de fundación por qué la gobernanza y la infraestructura posteriores resultaron tan complejas.

\subsection{Slashdot, Delicious y el ciclo mínimo de producto}

El ponente describe el producto original como una combinación de Slashdot y Delicious. El valor de Slashdot provenía de las discusiones de la comunidad técnica bajo cada noticia; Delicious, por su parte, mostraba la posibilidad de que los usuarios guardaran enlaces públicamente, los etiquetaran y los ordenaran en conjunto. Reddit eliminó a los editores tradicionales y dejó la selección del contenido en manos de los envíos y los votos, de modo que «quién decide la portada» dejó de ser un puesto centralizado y se volvió un conjunto de reglas de participación calculables. En ese momento el subreddit todavía no era el protagonista explícito del producto: las comunidades surgieron primero a partir de interacciones repetidas.

\begin{figure}[H]
\centering
\includegraphics[width=0.94\textwidth]{images/01-origin-product-loop.png}
\caption{El ciclo original de producto de Reddit: envío, votación, ordenamiento y discusión se impulsan mutuamente.}
\figsource{Redibujado a partir de los subtítulos de la entrevista 00:00--05:00, `images/01-origin-product-loop.png`.}
\end{figure}

\begin{knowledgebox}{Lectura de la figura: los cuatro estados del ciclo}
Submit genera contenido candidato, vote agrega las preferencias de los usuarios, rank decide cómo se reparte una atención limitada y discuss convierte un clic aislado en una relación continua. Los tres primeros pasos pueden resolverse con una base de datos y una lógica de ordenamiento; el cuarto da lugar a normas, reputación y conflictos. Cuando el sistema crece, el problema ya no es solo ordenar bien los enlaces, sino mantener las fronteras entre miles de espacios de discusión.
\end{knowledgebox}

\begin{importantbox}{Términos clave: UGC y subreddit}
UGC es user-generated content, es decir, contenido generado por los usuarios, que incluye publicaciones, comentarios, imágenes, votos y reglas de la comunidad. Un subreddit es la unidad de comunidad que se forma dentro de Reddit en torno a un tema, con sus propios moderadores, reglas y cultura de sus miembros. La plataforma proporciona la infraestructura global y los subreddits se encargan de gran parte de la gobernanza local; ambos niveles de poder se influyen mutuamente en los casos límite.
\end{importantbox}

\subsection{Cold start: cuándo pueden los fundadores salir del ciclo}

\term{Cold start} se refiere a la situación en que un sistema todavía no tiene suficientes usuarios, contenido o interacción, de modo que los usuarios nuevos no perciben su valor al entrar. En sus inicios, la portada de Reddit solo mostraba enlaces de las últimas 24 horas; si los fundadores no la alimentaban continuamente, la página podía quedar completamente vacía. Por eso Huffman y Alexis Ohanian usaban scripts para recopilar enlaces y luego los enviaban desde varias cuentas, para que los visitantes vieran un producto que ya estaba en funcionamiento. Esta práctica resolvía un problema de presentación y de aprendizaje, pero también planteaba un límite de autenticidad: el contenido semilla debe ceder su lugar a la participación real lo antes posible, y no puede usarse para simular de forma permanente el tamaño de la comunidad.

\begin{knowledgebox}{Nota de clase: quedarse dormido como criterio de crecimiento}
Huffman recuerda que un fin de semana no ejecutó el script diario de llenado y, al abrir el sitio, encontró la portada llena de enlaces enviados por nombres de usuario desconocidos. Ese momento fue más significativo que el número de registros: el sistema siguió funcionando sin el trabajo de ese día de los fundadores. En los mercados bilaterales, las comunidades y los productos de datos, el verdadero product-market fit suele manifestarse en que el ciclo central se sostiene por sí mismo.
\end{knowledgebox}

\begin{table}[H]
\centering
\begin{tabular}{L{0.2\textwidth}L{0.31\textwidth}L{0.39\textwidth}}
\toprule
Etapa & Acción de los fundadores & Señal de salida que debe observarse \\
\midrule
Periodo vacío & Preparar a mano o con scripts semillas de alta calidad & Los usuarios nuevos entienden para qué sirve el producto \\
Periodo de impulso & Crear varios temas y ejemplos de interacción & Los usuarios empiezan a imitar la conducta central \\
Periodo crítico & Reducir la frecuencia del llenado manual & La portada o la oferta se siguen actualizando sin supervisión \\
Periodo autónomo & Pasar a la prevención de abusos, las herramientas y el apoyo a la comunidad & La retención proviene del intercambio de valor entre usuarios \\
\bottomrule
\end{tabular}
\caption{El arranque en frío no es una campaña de crecimiento aislada, sino un proceso en el que el trabajo de los fundadores sale gradualmente del ciclo central.}
\end{table}

\subsection{Resumen de la sección}

La ventaja inicial de Reddit no fueron sus funciones complejas, sino un ciclo que los usuarios podían hacer suyo. Los fundadores pueden aportar temporalmente la oferta semilla, pero el criterio definitivo es si usuarios desconocidos siguen enviando, votando y discutiendo. Una vez formado el ciclo, la plataforma empieza a acumular estado de comunidad, y ese estado persiste a través de los múltiples cambios en la propiedad de la empresa y en su personal.
\section{El producto continúa, pero la organización se reinicia varias veces}

La base de datos y el dominio de una plataforma comunitaria pueden mantenerse sin interrupción, mientras que la organización responsable de ella puede pasar por una adquisición, la salida de su personal, rondas de financiamiento, crisis de liderazgo y una nueva expansión. Reddit pasó del programa de YC en 2005 a ser adquirida en 2006; el equipo fundador se fue alrededor de 2009; después, la plataforma se separó, obtuvo financiamiento y cambió de dirección, y Huffman regresó en medio de la crisis de 2015. Esta cronología muestra que el estado del sistema no reside solo en el código, sino también en la confianza de la comunidad, la experiencia de los empleados, las políticas pendientes y la deuda técnica acumulada a largo plazo.

\begin{figure}[H]
\centering
\includegraphics[width=0.95\textwidth]{images/02-company-timeline.png}
\caption{Etapas organizacionales de Reddit: fundación, adquisición, salida de los fundadores, financiamiento independiente, regreso en la crisis y empresa pública.}
\figsource{Redibujado a partir de los subtítulos de la entrevista, 05:00--10:00 y 37:30--39:30, `images/02-company-timeline.png`.}
\end{figure}

\begin{knowledgebox}{Lectura de la figura: separar la historia del producto de la historia de la organización}
Los usuarios ven un sitio que sigue disponible, pero internamente hubo varias reorganizaciones del control y de las capacidades. Después de 2009, el equipo llegó a ser muy pequeño y, aun así, la plataforma siguió creciendo gracias a la red comunitaria ya existente; cuando él regresó en 2015, la empresa ya tenía mucho más personal, pero carecía de un producto móvil maduro y de sistemas de crecimiento, datos, seguridad y gobernanza. Los efectos de red pueden retrasar el fracaso organizacional, pero no saldan por sí solos la deuda técnica ni la deuda de políticas.
\end{knowledgebox}

\subsection{Al regresar, primero diagnosticar por qué se teme cambiar}

Según Huffman, el obstáculo principal no era que los empleados desconocieran los problemas, sino que todos creían que Reddit era muy frágil y que cualquier cambio de políticas o de producto podía destruir el crecimiento. Este estado es común en los sistemas heredados de gran tamaño: la falta de observabilidad, de pruebas y de registros de decisiones impide que el equipo distinga entre los acoplamientos realmente peligrosos y los tabúes que se formaron con el tiempo. En esa situación, «mantener el statu quo» parece prudente, pero en realidad permite que los daños externos y la deuda interna sigan acumulándose.

\begin{importantbox}{Nota de clase: no cambiar también es una decisión de alto riesgo}
Al regresar, el ponente puso la política de contenido entre sus tareas principales, porque las comunidades extremistas ya habían sumido a la plataforma en una crisis permanente. El equipo temía que un cambio matara a Reddit; su diagnóstico, en cambio, era que el statu quo estaba desgastando la plataforma. Desde la ingeniería, este tipo de decisión debería exponer ambas curvas de riesgo a la vez: el cambio puede causar fallas a corto plazo, y no cambiar produce una pérdida segura a largo plazo.
\end{importantbox}

\begin{table}[H]
\centering
\begin{tabular}{L{0.2\textwidth}L{0.34\textwidth}L{0.36\textwidth}}
\toprule
Objeto de la reconstrucción & Capacidad que debe recuperarse & Señal verificable \\
\midrule
Producto & Móvil, experimentos, búsqueda y crecimiento & Frecuencia de lanzamientos, retención y reversión ante fallas \\
Infraestructura & Capacidad, recuperación ante fallas, conmutación de tráfico & Simulacros de picos, tiempo de recuperación, cobertura de degradación \\
Gobernanza & Reglas, herramientas, revisión entre equipos & Tiempo de atención, consistencia, resultado de las apelaciones \\
Organización & Responsables, rutas de escalamiento, registros de decisiones & Análisis posteriores a incidentes y calidad de la colaboración entre equipos \\
Relación con la comunidad & Comunicación transparente y herramientas para moderadores & Tasa de participación, encuestas de confianza, conflictos recurrentes \\
\bottomrule
\end{tabular}
\caption{Regresar no significa restaurar la autoridad del fundador, sino reconstruir capacidades organizacionales observables y modificables.}
\end{table}

\subsection{Poco capital no significa bajo costo}

Al principio, Reddit obtuvo unos 12,000 dólares de YC y después recibió un financiamiento pequeño y el apoyo de contratos comerciales. La baja inversión en efectivo suele idealizarse, pero el tiempo no remunerado de los fundadores, el contenido generado a mano, la paciencia de los primeros usuarios y el trabajo de seguridad pospuesto son costos reales. La adquisición dio al joven equipo efectivo y certeza de supervivencia, pero también introdujo el producto en la estructura organizacional de una empresa más grande. Al evaluar la eficiencia inicial, hay que volver a contabilizar los costos que se trasladaron a terceros o se aplazaron.

\begin{warningbox}{No derivar una fórmula de financiamiento a partir de un caso de supervivencia}
Que Reddit pudiera arrancar con muy poco dinero se relaciona con el costo de infraestructura de un producto web en 2005, las habilidades de los fundadores, la difusión que le dio Paul Graham y el tipo de producto. Hoy, para productos que involucran video, entrenamiento de modelos, pagos o datos sujetos a una regulación estricta, el costo mínimo viable es completamente distinto. La lección transferible es validar cuanto antes el ciclo central, no copiar mecánicamente un presupuesto de 12,000 dólares.
\end{warningbox}

\subsection{Resumen de la sección}

El ciclo de producto de Reddit atravesó varias transformaciones organizacionales, lo que muestra que una red comunitaria puede durar más que la estructura de la empresa. También muestra que el crecimiento no genera por sí solo gobernanza, seguridad ni capacidad de modificación. A partir de 2015, el trabajo central fue convertir una plataforma que sobrevivía por inercia en un sistema capaz de definir reglas claras, aplicarlas y soportar el cambio.
\section{De la autogestión por votación a la gobernanza por niveles}

Los primeros principios de Reddit se acercaban a «la plataforma no elimina; los usuarios deciden la visibilidad mediante votos». Cuando la escala era pequeña y la participación maliciosa limitada, esta autogestión reducía el poder editorial centralizado y además alentaba a las comunidades a formar su propia cultura. Al crecer la escala, la votación no podía atender el acoso entre comunidades, el contenido ilegal, la filtración de datos privados ni los grupos cuyo objetivo expreso es causar daño. La causa no es que los usuarios dejaran de ser confiables de repente, sino que el daño tiene externalidades: los afectados pueden no pertenecer a la comunidad que vota, y unos pocos coordinadores pueden manipular una mayoría local.

\subsection{Por qué la gobernanza tiene que organizarse por niveles}

Las reglas de la plataforma se encargan de la seguridad mínima, la legalidad y los límites entre comunidades; los moderadores de cada subreddit se encargan del tema y las normas locales; los votos de los usuarios se encargan del ordenamiento y la retroalimentación cotidianos. Cada uno de los tres niveles dispone de información distinta: la plataforma ve los patrones entre sitios y los riesgos legales, los moderadores entienden el contexto y la historia de la comunidad, y los usuarios comunes son quienes están más cerca del valor del contenido. Si cualquiera de los niveles concentra todo el poder, el resultado se distorsiona; el objetivo del sistema de gobernanza es que cada problema se atienda en el nivel que cuenta con información suficiente y cuyo alcance corresponde al del problema.

\begin{figure}[H]
\centering
\includegraphics[width=0.92\textwidth]{images/03-governance-stack.png}
\caption{Pila de gobernanza comunitaria: reparto de funciones entre la política de todo el sitio, la aplicación de normas por la plataforma, las reglas de cada subreddit y la participación de los usuarios.}
\figsource{Redibujado a partir de los subtítulos de la entrevista, 08:30--15:00, `images/03-governance-stack.png`.}
\end{figure}

\begin{knowledgebox}{Lectura de la figura: cuanto más alto el nivel, menos reglas pero límites más firmes}
La Site-wide policy se ocupa de los daños que atraviesan todo el sitio; debe ser relativamente estable, explicable y estar sujeta a la ley y a la responsabilidad pública. Las Subreddit rules pueden ser más detalladas y permiten que cada comunidad elija el formato de las publicaciones, los temas y el estilo de interacción. Voting es una señal débil y de alta frecuencia: sirve para ordenar, no para dictaminar sobre todos los problemas de seguridad. Las fallas de gobernanza suelen ocurrir por un desajuste de niveles: dejar que la votación resuelva la exposición de datos personales, o que la sede de la plataforma decida el tono de cada comunidad pequeña.
\end{knowledgebox}

\begin{importantbox}{Términos de primera aparición: moderation y doxing}
Moderation es la gobernanza del contenido y de la comunidad: incluye establecer reglas, revisar denuncias, restringir cuentas, atender apelaciones y apoyar a los moderadores. Doxing es publicar sin consentimiento información sensible de una persona, como su identidad, su domicilio o sus datos de contacto, lo que puede convertir un conflicto en línea en un daño en el mundo real. No puede decidirse solo por popularidad: requiere restricciones a nivel de todo el sitio y una respuesta rápida.
\end{importantbox}

\subsection{``Specifically vague'': por qué las reglas no pueden redactarse como una lista exhaustiva}

Los participantes maliciosos buscan resquicios en la letra de las reglas. Si una regla solo enumera palabras o conductas fijas, el atacante puede cambiar la ortografía, dividir sus acciones u organizar a otros para causar el mismo daño. Una regla demasiado amplia, en cambio, impide que los usuarios comunes anticipen dónde está el límite, y el equipo encargado de aplicarla tiende a actuar de forma inconsistente. Huffman usa ``specifically vague'' para describir el intervalo de trabajo: explicar con claridad el daño que se quiere evitar y el principio que lo sustenta, y conservar a la vez margen para juzgar, según el contexto, las conductas de evasión equivalentes.

\begin{figure}[H]
\centering
\includegraphics[width=0.92\textwidth]{images/04-policy-specificity.png}
\caption{El diseño de políticas necesita encontrar un intervalo de trabajo entre la previsibilidad y la resistencia a la evasión literal.}
\figsource{Redibujado a partir de los subtítulos de la entrevista, 13:40--16:00, `images/04-policy-specificity.png`.}
\end{figure}

\begin{knowledgebox}{Lectura de la figura: el margen de ambigüedad debe estar acotado por procedimientos}
La zona intermedia no autoriza a quien aplica las normas a decidir arbitrariamente. Necesita ejemplos, precedentes, revisión entre equipos, apelaciones, informes de transparencia y revisiones periódicas que la acoten. Cuanto más se dependa del juicio contextual, más hay que registrar por qué casos similares recibieron un trato distinto. El texto de la política expresa el objetivo; la justicia procedimental evita que el objetivo quede capturado por sesgos personales.
\end{knowledgebox}

\begin{warningbox}{Nota de clase: cada palabra proviene de un costo}
El ponente subraya que las frases de una política de contenido madura suelen ser el resultado condensado de años de incidentes, controversias y casos límite. Al leer diez reglas, no basta con evaluar si la redacción es elegante: también hay que preguntar si existen las herramientas correspondientes de detección, revisión humana, notificación, apelación y restablecimiento. Una regla sin un sistema que la aplique solo traslada la responsabilidad a los equipos de primera línea.
\end{warningbox}

\subsection{La aplicación de normas pasa de un interruptor binario a una escalera}

Las primeras herramientas se reducían a la suspensión permanente o a no hacer nada. Una elección binaria lleva al equipo a aplicar sanciones excesivas o, por temor a una pérdida irreversible, a no actuar en absoluto. Los sistemas más maduros forman una escalera con la explicación, la advertencia, la restricción de corto plazo, el subreddit timeout, la suspensión temporal y el retiro permanente. La intensidad de la intervención aumenta según el daño, el número de reincidencias, la disposición a cooperar y la posibilidad de recuperación, y se ofrece una vía de reversión para los errores de juicio.

\begin{figure}[H]
\centering
\includegraphics[width=0.94\textwidth]{images/05-enforcement-ladder.png}
\caption{Aplicación gradual de normas: de la explicación y la advertencia se escala progresivamente hasta la restricción temporal y el retiro permanente.}
\figsource{Redibujado a partir de los subtítulos de la entrevista, 16:00--20:00, `images/05-enforcement-ladder.png`.}
\end{figure}

\begin{importantbox}{Lectura de la figura: la precisión de las herramientas determina la precisión de la gobernanza}
Si el sistema solo cuenta con el ban permanente, el equipo de políticas no puede expresar diferencias aunque las entienda. El timeout permite a la plataforma pausar una comunidad que se salió de control, dialogar con los moderadores y conservar la posibilidad de restablecerla; la suspensión escalonada da a quien infringe de forma ocasional la oportunidad de corregirse y, al mismo tiempo, acumula evidencia sobre la conducta maliciosa persistente. Por ello, la calidad de la gobernanza resulta de la combinación de políticas, datos y herramientas de producto.
\end{importantbox}

\begin{table}[H]
\centering
\begin{tabular}{L{0.18\textwidth}L{0.27\textwidth}L{0.43\textwidth}}
\toprule
Intervención & Situación aplicable & Procedimiento que debe acompañarla \\
\midrule
Explicación/recordatorio & Primera vez, bajo riesgo, regla poco clara & Señalar la regla y el contenido concretos \\
Advertencia & Infracción confirmada pero corregible & Conservar el registro y explicar las consecuencias de reincidir \\
Restricción temporal & Infracción reiterada o riesgo de corto plazo & Precisar plazo, alcance y condiciones de restablecimiento \\
Timeout de la comunidad & Los moderadores perdieron el control o el conflicto escala & Dialogar con los moderadores, proteger el contenido histórico, revisar antes de reabrir \\
Retiro permanente & Daño grave, evasión maliciosa o fallas persistentes & Revisión de alto nivel, conservación de evidencia, apelación e informe de transparencia \\
\bottomrule
\end{tabular}
\caption{La escalera de aplicación de normas reúne en una misma tabla el daño, la reversibilidad y las garantías procedimentales.}
\end{table}

\subsection{Resumen de la sección}

La votación sirve para ordenar, pero no basta para asumir toda la responsabilidad sobre la seguridad y la legalidad. La evolución de la gobernanza de Reddit organizó por niveles la política de todo el sitio, la aplicación de normas por la plataforma, las reglas de cada comunidad y la participación de los usuarios, y recurrió a ``specifically vague'' y a herramientas graduales para atender la evasión maliciosa y los casos límite. La siguiente sección trata otro límite: si el hecho de que un contenido sea de lectura pública significa que cualquier entidad puede extraerlo sin límite y explotarlo comercialmente.
\section{Contenido público, licencias de datos y comunidades de pago}

Las plataformas abiertas suelen confundir «público en la web» con «uso ilimitado». Que una persona lea una publicación, que un buscador la indexe y devuelva tráfico, que un bot de moderación llame a la API, que un cliente comercial sustituya al producto oficial o que una empresa de modelos extraiga en masa datos de entrenamiento son actos con efectos muy distintos sobre los costos de la plataforma y sobre el intercambio de valor. Esta sección entiende la política de datos como un access contract: quién es el sujeto, qué escala tiene el acceso, si sustituye a la plataforma, si conserva el contexto y cómo asume los costos de infraestructura y de seguridad.

\subsection{Un mismo corpus requiere contratos de acceso distintos}

\term{Scraping} es la extracción automatizada del contenido de páginas web o de interfaces; \term{API} significa Application Programming Interface, es decir, la interfaz de llamadas acordada entre programas. Una API puede ofrecer identidad, permisos, límites de tasa y una estructura estable, mientras que la extracción de páginas web por lo general solo puede imitar la navegación de una persona. Que un contenido público pueda leerlo cualquier persona no implica que la plataforma deba soportar sin condiciones una extracción de cualquier escala, ni que quien extrae pueda ignorar las eliminaciones hechas por los autores, la privacidad y el contexto de la comunidad.

\begin{figure}[H]
\centering
\includegraphics[width=0.94\textwidth]{images/06-data-access-boundaries.png}
\caption{Cuando personas, buscadores, desarrolladores y entrenamiento de IA acceden al mismo corpus público, cada uno requiere una relación de intercambio distinta.}
\figsource{Redibujado a partir de los subtítulos de la entrevista 22:00--27:30, `images/06-data-access-boundaries.png`.}
\end{figure}

\begin{knowledgebox}{Lectura de la figura: la escala, la sustitución y el retorno definen los límites}
La navegación de una persona suele consumir recursos limitados y participar en la comunidad; los buscadores indexan en masa, pero aportan enlaces y tráfico; las herramientas para desarrolladores pueden potenciar a moderadores y usuarios, pero también pueden replicar un cliente completo; el entrenamiento de IA extrae texto a gran escala y traslada el valor al modelo. La política de acceso debe organizarse en niveles según el comportamiento real, y no reducirse a dicotomías como «empresa/individuo» o «abierto/cerrado».
\end{knowledgebox}

\begin{warningbox}{Público no significa perder el contexto}
Los comentarios de Reddit suelen contener discusiones, sarcasmo, información sensible al tiempo y normas internas de la comunidad. Al extraer un texto aislado, un modelo o un resumen de búsqueda puede presentar una opinión minoritaria como consenso. Quien usa los datos debe conservar la fuente, la fecha, la comunidad y el estado de eliminación; la plataforma, por su parte, debe explicar a los usuarios qué usos están autorizados y cómo se gestionan la exclusión voluntaria y las correcciones.
\end{warningbox}

\subsection{Subreddits de pago: ¿una herramienta adicional o una estratificación?}

La idea de comunidades de pago que se plantea en la entrevista no consiste en encerrar todos los subreddits públicos existentes tras un muro de pago, sino en dar una opción más a las comunidades dispuestas a ofrecer contenido especializado adicional o un servicio continuo. Los ingresos por suscripción pueden compensar el trabajo de moderadores, expertos y creadores, pero también pueden convertir la gobernanza en una protección de ingresos y hacer que a los críticos se les trate como usuarios que «afectan el negocio». El diseño debe centrarse en la creación voluntaria, un valor claro, los reembolsos, la comisión de la plataforma y la independencia de las reglas.

\begin{figure}[H]
\centering
\includegraphics[width=0.9\textwidth]{images/07-paid-community-loop.png}
\caption{El segundo ciclo de valor de las comunidades de pago: los miembros pagan, la comunidad entrega y los creadores y moderadores reciben una retribución.}
\figsource{Redibujado a partir de los subtítulos de la entrevista 27:50--32:30, `images/07-paid-community-loop.png`.}
\end{figure}

\begin{knowledgebox}{Lectura de la figura: los ingresos no compran excepciones a las reglas}
Tras pagar, los miembros esperan contenido y servicio estables; al recibir ingresos, los creadores aumentan su dedicación, y la plataforma aporta pagos, descubrimiento y capacidades de seguridad. Un ciclo sano exige que las políticas públicas sigan aplicándose y que los moderadores no relajen los límites frente al acoso o el fraude porque los usuarios de pago sean valiosos. También hay que evitar que la plataforma monetice de forma unilateral el conocimiento público que originalmente era gratuito y que acumularon voluntarios.
\end{knowledgebox}

\begin{importantbox}{Lista de verificación de diseño}
Una comunidad de pago debe responder, como mínimo: quién es dueño del contenido histórico; si los usuarios pueden exportarlo cuando la comunidad cierra; cómo se reparten los ingresos entre moderadores y creadores; cómo se gestionan los menores de edad y los consejos de alto riesgo; si el pago influye en el orden de las recomendaciones; si los usuarios pueden ver las vías de reembolso y de apelación. Sin estas interfaces, el botón de suscripción convierte la deuda de gobernanza en disputas financieras.
\end{importantbox}

\subsection{Resumen de la sección}

Una página web pública, una API abierta y una licencia comercial no son sinónimos. La plataforma necesita diseñar contratos distintos según la escala del acceso, el grado de sustitución, el costo de infraestructura y el retorno de valor. Las comunidades de pago tampoco son una simple función de ingresos: modifican el poder entre moderadores, miembros, creadores y la plataforma. El siguiente capítulo aplica estos principios al episodio más polémico: la fijación de precios de la API.
\section{Precios de la API: bien común del ecosistema y sustitución comercial}

La API de Reddit atiende a la vez a bots de moderadores, herramientas de investigación, clientes de accesibilidad, proyectos de aficionados, aplicaciones de terceros y usos comerciales de datos a gran escala. Un precio único parece sencillo, pero perjudica a las herramientas de bajo costo y alto valor público; que todo sea gratuito, en cambio, puede dejar a la plataforma con costos de infraestructura cada vez mayores y sin ingresos por publicidad o por licencias de datos. El verdadero problema de diseño es cómo clasificar, medir, notificar y migrar, no solo si «cobrar es correcto».

\subsection{Clasificar a los usuarios por su comportamiento, no por su identidad}

\term{Rate limit} es un límite al número de llamadas por unidad de tiempo, a la concurrencia o al volumen de datos, que sirve para proteger la capacidad y la equidad. El diseño de la API puede asignar un mayor peso de valor público a las herramientas de moderadores y a las aplicaciones de accesibilidad, ofrecer cuotas gratuitas a la investigación y a los proyectos de aficionados, y establecer contratos para los usos comerciales que sustituyen al cliente oficial o extraen datos de forma masiva. Los criterios de clasificación deben ser observables; de lo contrario, el equipo solo puede fijar los precios según la fama de cada empresa y su capacidad de negociación.

\begin{figure}[H]
\centering
\includegraphics[width=0.96\textwidth]{images/08-api-economics.png}
\caption{Niveles de la economía de la API: las herramientas comunitarias, la investigación, los clientes de terceros y el uso comercial a gran escala requieren contratos distintos.}
\figsource{Redibujada a partir de los subtítulos de la entrevista, 32:55--38:00, `images/08-api-economics.png`.}
\end{figure}

\begin{importantbox}{Lectura de la figura: antes del precio, definir la relación de sustitución}
Un moderation bot fortalece a Reddit mismo; un cliente completo de terceros puede sustituir la publicidad y la experiencia del producto oficial; el entrenamiento masivo de modelos traslada el valor del corpus a sistemas externos. Con el mismo número de llamadas, el significado económico puede ser distinto. Un esquema defendible debe hacer públicos las categorías, las cuotas gratuitas, las garantías de rendimiento, las solicitudes de excepción y los periodos de migración, en lugar de tratar a todos los usuarios con un precio unitario único en el último momento.
\end{importantbox}

\subsection{La protesta es retroalimentación, no un veto automático}

Los cambios en la API provocaron el cierre de subreddits y protestas de la comunidad. En la entrevista, Huffman reconoce que las protestas son coherentes con el carácter democrático de Reddit, aunque sostiene que la dirección sigue teniendo que tomar decisiones sobre la sostenibilidad de la plataforma. Ambas cosas pueden ser ciertas a la vez: la comunidad tiene derecho a expresar sus costos colectivos, y la plataforma también debe evaluar la seguridad, los ingresos, la infraestructura y el control del producto a largo plazo. Lo que importa es si el proceso de decisión permite que los afectados entiendan las razones con anticipación, prueben alternativas y dispongan de un plazo de migración razonable.

\begin{knowledgebox}{Nota de clase: no reducir el conflicto a un relato de buenos y malos}
Los desarrolladores de terceros invirtieron años en construir productos, los usuarios desarrollaron dependencia y las herramientas de moderadores asumen trabajo público de la plataforma; Reddit, por su parte, paga los costos de almacenamiento, ancho de banda, seguridad e ingeniería, y enfrenta la sustitución comercial. El conflicto surge de volver a ponerle precio a un antiguo contrato implícito. La revisión posterior más útil no es determinar quién está más enojado, sino identificar qué dependencias no se registraron, qué grupos no tenían una ruta de migración y qué costos nunca fueron transparentes.
\end{knowledgebox}

\begin{table}[H]
\centering
\begin{tabular}{L{0.2\textwidth}L{0.34\textwidth}L{0.36\textwidth}}
\toprule
Etapa & Lo que debe aportar la plataforma & Lo que debe aportar el ecosistema \\
\midrule
Descubrimiento & Clasificación de usos, modelo de costos, inventario de dependencias & Volumen real de llamadas, funciones críticas, necesidades de accesibilidad \\
Diseño & Precios, nivel gratuito, SLA, condiciones de excepción & Alternativas y periodo de migración aceptable \\
Piloto & Entorno aislado de pruebas, monitoreo, condiciones de reversión & Pruebas con poco tráfico y retroalimentación de compatibilidad \\
Migración & Documentación, herramientas, soporte y fechas definidas & Actualización de clientes, avisos a usuarios, exportación de datos \\
Revisión posterior & Interrupciones, quejas, ingresos y pérdidas del ecosistema & Casos fallidos, dependencias heredadas y propuestas de mejora \\
\bottomrule
\end{tabular}
\caption{Un cambio en la política de la API es un proyecto de migración de plataforma, no un anuncio de precios.}
\end{table}

\begin{warningbox}{Las restricciones posteriores a la salida a bolsa cambian el criterio de decisión}
Una empresa que cotiza en bolsa debe rendir cuentas de sus ingresos, costos y riesgos, pero la responsabilidad ante los accionistas no anula automáticamente la responsabilidad ante la comunidad. Los moderadores voluntarios, los creadores de contenido y los desarrolladores de largo plazo aportaron valor a la plataforma, aunque no necesariamente aparezcan en una tabla de costos tradicional. Una gestión madura debe incorporar la salud del ecosistema como indicador observable, en lugar de descubrir las dependencias solo cuando estalla la protesta.
\end{warningbox}

\subsection{Resumen de la sección}

La causa de fondo de la controversia de la API es que distintos patrones de uso compartían una interfaz que históricamente fue casi gratuita. Una solución sostenible requiere clasificación por comportamiento, medición transparente, excepciones por valor público, contratos comerciales y un procedimiento de migración. La apertura no significa precio cero para siempre, y la sostenibilidad tampoco equivale a un cierre repentino; entre ambas hace falta una gobernanza de la interfaz predecible a largo plazo.
\section{Picos de tráfico, Multi-Cloud y degradación elegante}

En las grandes plataformas comunitarias, los picos suelen originarse en noticias de última hora, eventos deportivos o políticos, o enlaces externos, y es difícil predecir del todo cuándo y dónde ocurrirán. El objetivo de la infraestructura no es que todas las funciones mantengan la misma calidad ante cualquier pico. Se trata de proteger las rutas centrales de lectura y escritura, la identidad y las capacidades de seguridad, y de degradar después, de manera progresiva, las funciones costosas según su valor para el negocio. La experiencia de Reddit se resume en cuatro pasos: absorb, divert, degrade y recover.

\subsection{Ruta de resiliencia en cuatro pasos}

El primer paso absorbe los picos ordinarios mediante cachés, colas y margen de capacidad elástica. El segundo evita las fallas locales al distribuir el tráfico entre varias regiones o varias nubes. El tercero, cuando la capacidad no alcanza, pasa al modo de solo lectura o desactiva las rutas de alto costo. El cuarto restablece la escritura, reproduce las colas y verifica el estado. \term{Graceful degradation}, o degradación elegante, significa que, ante falta de recursos o fallas parciales, el sistema conserva sus capacidades más importantes en lugar de caerse por completo.

\begin{figure}[H]
\centering
\includegraphics[width=0.96\textwidth]{images/09-traffic-resilience.png}
\caption{Ruta de resiliencia del tráfico en una plataforma comunitaria: absorción, desvío, degradación y recuperación.}
\figsource{Redibujada a partir de los subtítulos de la entrevista, 40:00--43:20, `images/09-traffic-resilience.png`.}
\end{figure}

\begin{knowledgebox}{Lectura de la figura: el orden de degradación debe codificarse de antemano}
Si se discute qué funciones importan cuando el incidente ya está en curso, la presión y la información parcial dominan las decisiones. El sistema debe definir de antemano la prioridad de la ruta central de lectura, la ruta de escritura de publicaciones y comentarios, el inicio de sesión, la moderación de seguridad, la búsqueda, las recomendaciones y las notificaciones. El modo de solo lectura mantiene visible el contenido de la comunidad, pero también puede impedir que los moderadores atiendan un daño que se está propagando. Por eso las herramientas de seguridad no deben desactivarse en bloque junto con las funciones de escritura ordinarias.
\end{knowledgebox}

\subsection{Multi-cloud no es replicar un conjunto de máquinas}

Usar varias nubes reduce el riesgo de que falle un solo proveedor y el riesgo en la negociación con él. A cambio, aumenta el costo de la replicación de datos, la identidad, la red, la observabilidad y los simulacros. Una arquitectura multinube solo aporta capacidad de recuperación si el equipo puede redirigir el tráfico en un tiempo aceptable, confirmar la consistencia de los datos y hacer que los servicios críticos funcionen de manera independiente en el segundo entorno. Desplegar unos cuantos recursos en otra nube sin haber hecho nunca un simulacro se parece más a diversificar las compras que a una conmutación por error.

\begin{importantbox}{Nota de clase: los objetivos de recuperación deben expresarse como acciones del usuario}
«La segunda nube está disponible» no es lo bastante concreto. Un mejor objetivo sería este: los usuarios pueden leer las publicaciones populares, los moderadores pueden ocultar contenido peligroso, las sesiones iniciadas no se invalidan de forma masiva y los comentarios nuevos pueden ponerse en cola y escribirse después de la recuperación. Expresar el objetivo como acciones del usuario pone al descubierto las dependencias y también permite saber si los simulacros cubren de verdad el valor central.
\end{importantbox}

\begin{table}[H]
\centering
\begin{tabular}{L{0.19\textwidth}L{0.31\textwidth}L{0.4\textwidth}}
\toprule
Capacidad & Impacto en el usuario si falla & Estrategia de degradación \\
\midrule
Lectura de contenido & La plataforma pierde casi toda su utilidad & Caché en varias capas, instantáneas estáticas, solo lectura entre regiones \\
Publicaciones y comentarios & Se interrumpe la participación en tiempo real & Colas, limitación de tasa, escritura diferida \\
Herramientas de seguridad para moderadores & El daño no puede contenerse a tiempo & Capacidad dedicada reservada e interfaz mínima de moderación \\
Ordenamiento de recomendaciones & Disminuye la personalización & Volver a lo más popular de la comunidad o al orden cronológico \\
Búsqueda/resúmenes con IA & El descubrimiento de información se vuelve más lento & Desactivar los resúmenes y conservar la búsqueda por palabras clave \\
Notificaciones y análisis & Se admite cierto retraso & Procesamiento por lotes y reproducción después de la recuperación \\
\bottomrule
\end{tabular}
\caption{La degradación se diseña según el valor para el usuario, no según qué equipo es dueño de cada microservicio.}
\end{table}

\begin{warningbox}{La reparación de datos después de la recuperación también es importante}
Que el tráfico se haya redirigido con éxito no significa que el incidente haya terminado. Las colas pueden consumirse más de una vez, los contadores pueden desviarse, el orden de los votos y los árboles de comentarios pueden quedar inconsistentes por un tiempo, y los eventos de seguridad pueden acumularse durante la ventana de degradación. El proceso de recuperación requiere reconciliation, eliminación de duplicados, auditoría y compensación a los usuarios. No basta con ver que vuelve a subir la tasa de éxito de HTTP.
\end{warningbox}

\subsection{Resumen de la sección}

La resiliencia de una plataforma depende de la capacidad reservada, de una redirección del tráfico que pueda ensayarse en simulacros, de una degradación elegante ordenada por el valor para el usuario y de la verificación de los datos después de la recuperación. Multi-cloud solo aporta resiliencia cuando la aplicación, los datos y la identidad pueden cambiar de entorno de verdad. A continuación se analiza cómo la IA cambia los puntos de entrada del tráfico y la forma de las interfaces, y qué necesidades de la comunidad no podrán sustituirse con modelos de resumen.
\section{La era de la IA: Scrollers, Seekers y Agents}

La pregunta «¿la IA reemplazará a Reddit?» es demasiado general. La entrevista divide a los usuarios en scrollers y seekers: los primeros entran a la comunidad para participar, observar a otros y crear sentido de pertenencia; los segundos llegan con una pregunta concreta y buscan experiencias recientes. La búsqueda generativa cambia primero el camino del seeker, porque un resumen puede condensar varias publicaciones; difícilmente sustituye la motivación del scroller de interactuar con personas reales. La estrategia de producto debe partir de la intención del usuario, en lugar de tratar todas las vistas de página como un mismo tipo de tráfico.

\subsection{Dos contratos de usuario}

Los scrollers aceptan una navegación abierta, y su valor proviene de la novedad, la conversación y la identidad; los seekers quieren obtener respuestas rápido, pero siguen valorando la experiencia, los desacuerdos y la actualidad de los comentarios. Durante mucho tiempo, los motores de búsqueda llevaron a los seekers a Reddit; los resúmenes de IA pueden hacer fuera del sitio una parte de la lectura, pero también pueden reducir la confianza porque carecen de fuentes y de la estructura de la controversia. La defensa de Reddit no consiste solo en impedir la extracción automática de contenido, sino también en lograr que la búsqueda dentro del sitio exprese mejor el conocimiento de la comunidad.

\begin{figure}[H]
\centering
\includegraphics[width=0.92\textwidth]{images/10-scrollers-seekers.png}
\caption{Dos modos de uso: los scrollers que participan en la comunidad y los seekers que buscan respuestas.}
\figsource{Redibujada a partir de los subtítulos de la entrevista, 43:30--45:10, `images/10-scrollers-seekers.png`.}
\end{figure}

\begin{knowledgebox}{Lectura de la figura: una misma página puede atender intenciones distintas}
Una persona puede encontrar primero una recomendación de compra mediante una búsqueda y, después, suscribirse a la comunidad y convertirse en participante de largo plazo. El camino de búsqueda (seek) exige relevancia, evidencia y actualidad; el camino de exploración (scroll) exige identidad comunitaria, descubrimiento e interacción segura. Si el producto solo optimiza la tasa de clics de los resúmenes, puede dañar la oferta de discusión; si solo optimiza el tiempo en el feed, los usuarios que llegan con una pregunta no podrán localizar rápido respuestas de alta calidad.
\end{knowledgebox}

\subsection{Reddit Answers: primero recuperar, luego resumir y después enlazar de vuelta}

\term{RAG} es Retrieval-Augmented Generation, generación aumentada por recuperación. El sistema primero recupera fragmentos relevantes de un corpus externo y luego hace que el modelo de lenguaje genere una respuesta basada en esos fragmentos. Según la entrevista, Reddit Answers busca entre una gran cantidad de publicaciones y comentarios, genera un resumen integrado y enlaza cada afirmación factual con el comentario específico del que proviene. Estos enlaces de vuelta permiten al usuario revisar las fuentes, leer las objeciones y entender el contexto temporal; son un diseño importante cuando el corpus de una comunidad se incorpora a un producto de IA.

\begin{figure}[H]
\centering
\includegraphics[width=0.94\textwidth]{images/11-reddit-answers-rag.png}
\caption{El flujo básico de Reddit Answers: pregunta, recuperación, resumen y enlaces a las fuentes a nivel de comentario.}
\figsource{Redibujada a partir de los subtítulos de la entrevista, 45:00--46:10, `images/11-reddit-answers-rag.png`.}
\end{figure}

\begin{importantbox}{Lectura de la figura: el grounding no garantiza la verdad}
La recuperación puede impedir que el modelo invente contenido desvinculado del corpus, y los enlaces a comentarios aumentan la capacidad de auditoría; aun así, el contenido de la comunidad puede ser erróneo, estar desactualizado, haber sido manipulado o representar solo a una minoría. El sistema debe mostrar la fecha, la comunidad, la cobertura de la muestra y los desacuerdos, y no debe presentar «la tendencia mayoritaria de los comentarios» como una conclusión profesional. Las preguntas médicas, legales y de seguridad requieren además fuentes más sólidas y advertencias de riesgo.
\end{importantbox}

\begin{knowledgebox}{Nota de clase: un prototipo de 90 días demuestra que la interfaz era clara, no que el problema esté resuelto}
El ponente mencionó que Reddit Answers tuvo una versión rápida en unos 90 días. El índice de búsqueda, el corpus de contenido, el sistema de identidad y los puntos de entrada del producto ya existentes redujeron el costo del prototipo. El trabajo de largo plazo consiste en evaluar la calidad de los resúmenes, manejar contenido malicioso, respetar las eliminaciones, controlar la latencia y el costo, y evitar que las páginas de respuestas terminen reduciendo las contribuciones a la comunidad.
\end{knowledgebox}

\subsection{Agent API: el contrato de acceso de la siguiente generación}

Un agent no solo lee páginas: también compara, planea y ejecuta acciones en nombre del usuario. Hoy los agents suelen interpretar los sitios web a través de un navegador, lo cual es una interfaz accidental: su estructura es inestable, la expresión de permisos es burda, la plataforma difícilmente identifica el propósito y también le cuesta cobrar el costo. Una \term{Agent API} debería ofrecer identidad de máquina, autorización del usuario, acciones invocables, límites de tasa, atribución de fuentes, mecanismos de pago y de revocación, para que la plataforma sepa quién hace qué en nombre de quién.

\begin{figure}[H]
\centering
\includegraphics[width=0.94\textwidth]{images/12-agent-access-contract.png}
\caption{El contrato de acceso para agents debe alinear identidad, costo y valor entre la plataforma, la API dedicada y el agente que ejecuta.}
\figsource{Redibujada a partir de los subtítulos de la entrevista, 46:10--48:22, `images/12-agent-access-contract.png`.}
\end{figure}

\begin{knowledgebox}{Lectura de la figura: una interfaz para agents incluye al menos cuatro registros contables}
El registro de identidad indica quiénes son el agent, el usuario final y el desarrollador; el registro de permisos indica qué acciones se pueden leer, escribir y ejecutar; el registro de costos anota las invocaciones, el ancho de banda y el uso de datos de alto valor; el registro de atribución permite que las respuestas remitan a la comunidad y a los autores. Si el agent omite la publicidad o la interfaz oficial, la plataforma también debe buscar nuevas formas de cobro que no induzcan promoción encubierta.
\end{knowledgebox}

\begin{warningbox}{Nota de clase: la plataforma no tiene la opción de «rechazar la era de los agents»}
La valoración verbal de Huffman es que los agents ya llegaron y que las plataformas solo pueden elegir cómo adaptarse. Él sigue prefiriendo la apertura y la accesibilidad, pero reconoce que el acceso gratuito ilimitado puede causar problemas de infraestructura o de modelo de negocio. Una buena estrategia no consiste en declarar apertura total o cierre total, sino en establecer pronto interfaces dedicadas, costos claros y reglas de protección de la comunidad.
\end{warningbox}

\subsection{Resumen de la sección}

La IA cambia primero el camino de búsqueda y resumen de los seekers, mientras que la necesidad de los scrollers de una comunidad humana sigue existiendo. Reddit Answers muestra la combinación de recuperación, resumen y enlaces de vuelta a las fuentes; la agent API, por su parte, lleva la cuestión del acceso hacia la identidad, la autorización, las acciones y el cobro. El activo central de la plataforma no es solo el corpus de texto, sino la producción continua de experiencia humana con contexto temporal, comunitario y de interacción.
\section{Síntesis y ampliación}

Los veinte años de evolución de Reddit pueden verse como tres curvas de sistema entrelazadas. La curva de producto se extendió de la clasificación de enlaces a las comunidades, la búsqueda y la IA; la curva de gobernanza se extendió de «no eliminar» a reglas por niveles, aplicación graduada de sanciones y procedimientos de apelación; la curva de infraestructura se extendió de un solo sitio web a la multinube, la degradación del tráfico, la API y las interfaces para agents. Cada extensión reinterpreta qué significa «impulsado por los usuarios»: los usuarios disponen de más capacidad de expresión y de más herramientas, y al mismo tiempo la plataforma asume responsabilidades más explícitas de seguridad, confiabilidad y sostenibilidad.

\begin{importantbox}{Siete conclusiones transferibles}
\begin{enumerate}
  \item El punto final del arranque en frío es que el ciclo central siga funcionando después de que los fundadores se retiren, no que el número de registros alcance una cifra redonda.
  \item Los efectos de red pueden ocultar la deuda organizacional y técnica; el crecimiento sostenido no debe tomarse como prueba de la salud del sistema.
  \item El voto es una señal de clasificación, no un mecanismo de gobernanza universal; las externalidades requieren reglas de nivel superior.
  \item El texto de las políticas, las herramientas de aplicación, las apelaciones y la transparencia determinan en conjunto la calidad de la moderation.
  \item La lectura pública, la API gratuita y la extracción comercial ilimitada son tres contratos de acceso distintos.
  \item El diseño para la resiliencia debe priorizar según las acciones de los usuarios e incorporar a la ruta central la seguridad y la verificación posterior a la recuperación.
  \item Los resúmenes de IA y los agents deben conservar la fuente, la autorización y el retorno de valor; de lo contrario, agotarán la aportación de la comunidad.
\end{enumerate}
\end{importantbox}

\subsection{Tarea práctica: diseñar un plan de cambio para una plataforma comunitaria}

Elige un foro real, una comunidad de código abierto o una plataforma de contenidos, y propón un cambio que modifique el contrato con los usuarios, por ejemplo, cobrar por la API, comunidades de pago, resúmenes de IA o escritura por parte de agents. La tarea debe abarcar a la vez el producto, la gobernanza, la infraestructura y el procedimiento de migración; no basta con presentar una tabla de beneficios comerciales. El lector debe poder identificar en tu plan los grupos afectados, los riesgos, las condiciones de reversión y los indicadores de éxito.

\begin{enumerate}
  \item Dibuja el ciclo central actual y señala qué participantes aportan el contenido, la revisión, la distribución y los ingresos.
  \item Define las categorías de acceso, la identidad, los permisos, las tasas y las reglas de liquidación posteriores al cambio.
  \item Enumera al menos tres modos de falla: uno debe provenir de un cambio en los incentivos y otro de una degradación de la infraestructura.
  \item Diseña los procesos de piloto, notificación, migración, apelación y reversión, y explica cómo se hará pública la revisión retrospectiva.
\end{enumerate}

\begin{knowledgebox}{Criterios de aceptación}
Un plan de alta calidad debe convertir «apertura», «comunidad», «equidad» y «sostenibilidad» en interfaces e indicadores observables. Si el plan solo dice «escuchar a los usuarios» sin explicar cómo la retroalimentación llega a las decisiones, o solo dice «recuperar costos» sin un modelo de costos ni reglas de excepción, los sistemas clave siguen ocultos detrás de consignas.
\end{knowledgebox}

\subsection{Lecturas adicionales}

El estudio puede continuar por tres vías. Primera: contrasta la página del curso de Stanford Winter 2025 con los subtítulos locales para construir una línea de tiempo de las decisiones clave de Reddit, y distingue con claridad qué proviene de los recuerdos del ponente y qué requiere verificación con fuentes externas. Segunda: lee las políticas de identidad, rate limit, SLA y retiro de una API pública, y evalúa si respaldan un ecosistema a largo plazo. Tercera: elige un producto de RAG o de agents y revisa si las respuestas enlazan a sus fuentes, si los usuarios pueden revocar la autorización y cómo maneja la plataforma la caché y los datos derivados después de una eliminación.

\section{Apéndice: formulario de revisión de cambios en plataformas comunitarias}

Este apéndice condensa los casos de esta clase en un proceso de revisión reutilizable, aplicable a cambios de alto impacto como políticas de contenido, API, funciones de pago, resúmenes generados por IA e interfaces para agent. La revisión no busca que todos los grupos lleguen a un acuerdo, sino exigir que el equipo haga públicos, antes del lanzamiento, las dependencias clave, los costos, los cambios de poder y las rutas de recuperación ante fallas. Cada punto debe tener un responsable, evidencia y condiciones de detención, para no tomar «la comunidad se adaptará» como supuesto por defecto.

\begin{table}[H]
\centering
\begin{tabular}{L{0.18\textwidth}L{0.33\textwidth}L{0.39\textwidth}}
\toprule
Dimensión de revisión & Pregunta obligatoria & Evidencia mínima \\
\midrule
Contrato con el usuario & Quién gana o pierde qué capacidades & Recorrido del usuario, diferencias de permisos, ejemplos de notificación \\
Poder en la comunidad & Cómo cambia el poder de moderadores, creadores, miembros y plataforma & Matriz de decisiones, rutas de apelación y de salida \\
Modelo económico & Quién asume los costos y quién obtiene los ingresos o el valor de los datos & Rangos de costo, nivel gratuito, reparto de ingresos y excepciones \\
Gobernanza de seguridad & Qué abusos y externalidades amplificará la nueva función & Modelo de amenazas, herramientas de moderación, proceso de escalamiento \\
Infraestructura & Qué acciones se conservan en picos de carga, fallas y reversiones & Pruebas de capacidad, orden de degradación, simulacros de recuperación \\
Ciclo de vida de los datos & Cómo se tratan los datos derivados tras la eliminación de contenido o la revocación de autorizaciones & Plazos de retención, limpieza de caché, registros de auditoría \\
\bottomrule
\end{tabular}
\caption{Todo cambio de alto impacto en una plataforma debe aprobar a la vez las revisiones de producto, gobernanza, economía, infraestructura y datos.}
\end{table}

\begin{table}[H]
\centering
\begin{tabular}{L{0.2\textwidth}L{0.32\textwidth}L{0.38\textwidth}}
\toprule
Etapa de lanzamiento & Indicadores adelantados & Condiciones de detención o reversión \\
\midrule
Pruebas internas & Tasa de permisos correctos, bloqueos erróneos, latencia y costo & Ampliación inexplicable de permisos o filtración de datos \\
Piloto en comunidades pequeñas & Adopción, quejas, horas de trabajo de los moderadores, tasa de apelaciones exitosas & Daño concentrado en grupos vulnerables o herramientas insuficientes \\
Ampliación del alcance & Tasa de migración del ecosistema, tasa de fallas, ingresos y oferta de contenido & Caída de la oferta principal, interrupción de clientes clave \\
Operación estable & Retención a largo plazo, coherencia de la gobernanza, recuperación de costos & Deuda en crecimiento constante sin un plan de corrección creíble \\
Recuperación de incidentes & Tiempo de recuperación, rezago acumulado, compensación y transparencia & Datos imposibles de conciliar o daños repetidos sin control \\
\bottomrule
\end{tabular}
\caption{Cada indicador debe ir acompañado de una condición de detención; de lo contrario, el monitoreo solo registrará las fallas sin cambiar las decisiones.}
\end{table}

\begin{warningbox}{Los tres grupos que con más facilidad se omiten en la revisión}
El primero son los desarrolladores que mantienen herramientas públicas sin tener un contrato comercial; el segundo, los moderadores voluntarios que asumen gran parte de la moderación sin figurar en la plantilla de personal; el tercero, los autores de contenido afectados por el entrenamiento con datos o por los resúmenes generados por IA, que ya no visitan directamente la plataforma. Si un cambio solo contabiliza a los clientes de pago y a los usuarios activos en las páginas, subestimará de forma sistemática las contribuciones y las pérdidas de estos grupos.
\end{warningbox}

\end{document}
